%% ma: L12-B08-0174
%% lop: 12
%% chuong: 2
%% bai: 8
%% dang: 12.8.53
%% loai: TN4
%% muc: VDC
%% dap_an: "B"
%% nguon: "CĐ2.4.pdf giáo viên gửi 27/10/2026 (M07, Lê Minh Khôi – Nguyễn Anh Khôi, Chương II Chủ đề 4 Các bài toán về tâm tỉ cự), Câu 24"
%% kien_thuc: "Bài 8 (khoảng cách hai điểm, biểu thức tọa độ), Bài 6 (bình phương vô hướng, tâm tỉ cự), lớp 10 (giá trị lớn nhất của tam thức bậc hai)"
%% trang_thai: cho-duyet
%% ly_do: "Dạng mới 12.8.53 đề xuất, chờ giáo viên duyệt"
%% ghi_chu: "Phương án gốc ghi nhầm ký hiệu T_max, đã sửa thành M_max"
\begin{ex}
Trong không gian $Oxyz$, cho ba điểm $A(2;4;0)$, $B(2;0;4)$, $C(2;-2;6)$. Tìm giá trị lớn nhất của biểu thức $M=50TA^2+150TB^2-\left(TA^2+TC^2\right)^2$ với $T$ là điểm bất kỳ trong không gian.
\choice{$M_{\max}=95$}{\True $M_{\max}=100$}{$M_{\max}=25$}{$M_{\max}=50$}
\loigiai{Gọi $J(2;1;3)$ là trung điểm $AC$; $J$ cũng thỏa $50\overrightarrow{JA}+150\overrightarrow{JB}=\vec 0$.

$TA^2+TC^2=2TJ^2+36$, $50TA^2+150TB^2=200TJ^2+50\cdot18+150\cdot2=200TJ^2+1200$.

Đặt $u=TJ^2\ge0$: $M=200u+1200-(2u+36)^2=-4u^2+56u-96\le100$ (khi $u=7$).}
\end{ex}
